\documentclass[a4paper,11pt]{article}
\usepackage[utf8]{inputenc}
\usepackage[T1]{fontenc}
\usepackage[french]{babel}
\usepackage{amsmath, amssymb, amsthm}
\usepackage{graphicx}
\usepackage{tcolorbox}
\usepackage{tikz}
\usetikzlibrary{shapes,arrows,positioning,calc}
\usepackage{listings}
\usepackage{xcolor}
\usepackage{geometry}
\usepackage{hyperref}

\geometry{hmargin=2.5cm,vmargin=2.5cm}

% Configuration des boîtes
\tcbuselibrary{skins,breakable}

\newtcolorbox{conceptbox}[1][]{
  colback=blue!5!white,
  colframe=blue!75!black,
  title={\textbf{💡 Concept Clé}},
  fonttitle=\bfseries,
  enhanced,
  drop shadow,
  breakable,
  #1
}

\newtcolorbox{mathbox}[1][]{
  colback=green!5!white,
  colframe=green!60!black,
  title={\textbf{📐 Fondement Mathématique}},
  fonttitle=\bfseries,
  enhanced,
  drop shadow,
  breakable,
  #1
}

\newtcolorbox{codebox}[1][]{
  colback=gray!5!white,
  colframe=gray!75!black,
  title={\textbf{💻 Implémentation Python}},
  fonttitle=\bfseries,
  breakable,
  #1
}

\newtcolorbox{warningbox}[1][]{
  colback=red!5!white,
  colframe=red!75!black,
  title={\textbf{⚠️ Attention}},
  fonttitle=\bfseries,
  enhanced,
  #1
}

% Configuration du code
\definecolor{codegreen}{rgb}{0,0.6,0}
\definecolor{codegray}{rgb}{0.5,0.5,0.5}
\definecolor{codepurple}{rgb}{0.58,0,0.82}
\definecolor{backcolour}{rgb}{0.95,0.95,0.92}

\lstdefinestyle{pythonstyle}{
    language=Python,
    backgroundcolor=\color{backcolour},   
    commentstyle=\color{codegreen},
    keywordstyle=\color{magenta},
    numberstyle=\tiny\color{codegray},
    stringstyle=\color{codepurple},
    basicstyle=\ttfamily\footnotesize,
    breakatwhitespace=false,         
    breaklines=true,                 
    captionpos=b,                    
    keepspaces=true,                 
    numbers=left,                    
    numbersep=5pt,                  
    showspaces=false,                
    showstringspaces=false,
    showtabs=false,                  
    tabsize=2
}
\lstset{style=pythonstyle}

\title{\textbf{Local Outlier Factor (LOF)} \\ \large Détection d'Anomalies basée sur la Densité Locale}
\author{Cours de Machine Learning Non Supervisé}
\date{\today}

\begin{document}

\maketitle
\tableofcontents
\newpage

\section{Introduction}

Le \textbf{Local Outlier Factor} (LOF) (Breunig et al., 2000) est une méthode pionnière de détection d'anomalies. Contrairement aux méthodes globaux (qui comparent un point à tous les autres), LOF s'intéresse à la \textbf{densité locale}.

\begin{conceptbox}
\textbf{Intuition Fondamentale :}

Une anomalie est un point qui est \textbf{plus isolé} que ses voisins.

\begin{itemize}
    \item Si un point a une densité similaire à celle de ses voisins, c'est un point normal.
    \item Si sa densité est \textbf{beaucoup plus faible} que celle de ses voisins, c'est une anomalie.
\end{itemize}

L'avantage majeur est de pouvoir détecter des anomalies dans des zones de densités différentes, là où un seuil de distance global échouerait.
\end{conceptbox}

\section{Principe de Fonctionnement}

LOF calcule un score pour chaque point, basé sur le rapport entre sa densité locale et la densité locale moyenne de ses $k$ plus proches voisins.

\subsection{Illustration}

\begin{center}
\begin{tikzpicture}[scale=1]
    % Cluster dense
    \foreach \i in {1,...,15}
      \fill[blue!70] ({random(0,2)+0.5}, {random(0,2)+0.5}) circle (2pt);
    \node at (1.5,1.5) [blue] {Cluster dense $C_1$};

    % Cluster épars
    \foreach \i in {1,...,10}
      \fill[blue!70] ({random(4,6)+0.5}, {random(4,6)+0.5}) circle (2pt);
    \node at (5.5,5.5) [blue] {Cluster épars $C_2$};
    
    % Anomalie locale près de C1
    \fill[red] (3.5, 2.5) circle (3pt) node[right, red] {$O_1$};
    
    % Anomalie globale
    \fill[red] (6, 1) circle (3pt) node[right, red] {$O_2$};
\end{tikzpicture}

\textit{Figure 1 : LOF permet de détecter $O_1$ comme anomalie relative à $C_1$, même si la distance entre points dans $C_2$ est similaire.}
\end{center}

\section{Fondements Mathématiques}

L'algorithme se déroule en 4 étapes clés :

\subsection{1. k-Distance et Voisins}
Pour un point $p$ et un entier $k$, la $k$-distance($p$) est la distance au $k$-ième voisin le plus proche.
$N_k(p)$ est l'ensemble de ces $k$ voisins (peut être $>k$ en cas d'égalité).

\subsection{2. Reachability Distance (Distance d'Atteignabilité)}
C'est une distance "lissée". Pour éviter les problèmes si des points sont superposés :
$$ \text{reach-dist}_k(p, o) = \max\{ \text{k-distance}(o), d(p, o) \} $$
où $o \in N_k(p)$.
Intuition : Si $p$ est très proche de $o$, on utilise la k-distance de $o$ comme "bulle de protection" minimale.

\subsection{3. Local Reachability Density (LRD - Densité Locale)}
La densité locale de $p$ est l'inverse de la distance d'atteignabilité moyenne de $p$ par rapport à ses voisins :
$$ \text{lrd}_k(p) = \frac{|N_k(p)|}{\sum_{o \in N_k(p)} \text{reach-dist}_k(p, o)} $$

\subsection{4. Local Outlier Factor (LOF)}
Le score final compare la LRD de $p$ à celles de ses voisins $o$ :
$$ \text{LOF}_k(p) = \frac{\sum_{o \in N_k(p)} \frac{\text{lrd}_k(o)}{\text{lrd}_k(p)}}{|N_k(p)|} = \frac{\text{Moyenne LRD voisins}}{\text{LRD du point}} $$

\begin{mathbox}
\textbf{Interprétation du Score LOF}

\begin{itemize}
    \item $\text{LOF} \approx 1$ : Densité similaire aux voisins $\Rightarrow$ \textbf{Normal}.
    \item $\text{LOF} \ll 1$ : Point plus dense que ses voisins $\Rightarrow$ \textbf{Intérieur de cluster}.
    \item $\text{LOF} \gg 1$ : Densité beaucoup plus faible que les voisins $\Rightarrow$ \textbf{Anomalie}.
\end{itemize}
En pratique, un seuil $> 1.5$ ou $> 2.0$ est souvent utilisé pour flagger une anomalie.
\end{mathbox}

\section{Hyperparamètre Principal : n\_neighbors}

Le seul paramètre critique est $k$ (\texttt{n\_neighbors} dans Scikit-Learn).
\begin{itemize}
    \item $k$ trop petit (ex: 1 ou 2) : Sensible au bruit, faux positifs locaux.
    \item $k$ trop grand : L'effet local se dilue, on tend vers une approche globale.
    \item \textbf{Recommandation} : Souvent $k=20$ par défaut. Breunig suggère une plage (minPts, maxPts).
\end{itemize}

\section{Comparaison avec Isolation Forest et SVM}

\begin{center}
\begin{tabular}{|l|l|l|}
\hline
\textbf{Méthode} & \textbf{Force} & \textbf{Faiblesse} \\
\hline
Isolation Forest & Rapide, Scalable, Global & Rate les anomalies locales denses \\
One-Class SVM & Frontières complexes & Lent, Paramétrage difficile \\
\textbf{LOF} & \textbf{Détecte anomalies locales} & Lent ($O(n^2)$), Pas de prédiction \\
\hline
\end{tabular}
\end{center}

\begin{warningbox}
\textbf{Nouveauté et Prédiction :}
Standard LOF est conçu pour la détection d'outliers dans un dataset statique (\textit{transductif}). Il ne prédit pas nativement sur de nouvelles données sans recalculer les distances.
$\Rightarrow$ Scikit-Learn propose \texttt{novelty=True} pour l'utiliser comme un estimateur qui peut appeler \texttt{predict()}, mais attention : cela change légèrement la logique.
\end{warningbox}

\section{Exemple de Code Python}

\begin{codebox}
\begin{lstlisting}
from sklearn.neighbors import LocalOutlierFactor
from sklearn.preprocessing import StandardScaler
from sklearn.metrics import classification_report, confusion_matrix
import numpy as np

# 1. PRÉPARATION
# Les méthodes de distance nécessitent une normalisation !
scaler = StandardScaler()
X_scaled = scaler.fit_transform(X)

# 2. MODÉLISATION
# n_neighbors=20 : standard
# contamination=0.03 : on s'attend à 3% d'anomalies
# novelty=False : On detecte sur le jeu d'entrainement lui-même (classique)
lof = LocalOutlierFactor(
    n_neighbors=20, 
    contamination=0.03,
    n_jobs=-1
)

print("Calcul du LOF...")
# fit_predict existe pour le mode outlier detection
preds = lof.fit_predict(X_scaled) 

# 3. CONVERSION
# LOF : 1 = Normal, -1 = Anomalie
y_pred_lof = np.where(preds == -1, 1, 0)

# Extraire les scores bruts (negative_outlier_factor_)
# Plus c'est négatif, plus c'est une anomalie (dans sklearn implementation)
lof_scores = -lof.negative_outlier_factor_ 

# 4. ÉVALUATION
print("Matrice de Confusion :")
print(confusion_matrix(y, y_pred_lof))

print("\nRapport :")
print(classification_report(y, y_pred_lof, target_names=['Normal', 'Panne']))

# Note : lof_scores > 1.5 indiquent généralement une anomalie
print(f"Score LOF max : {lof_scores.max():.2f}")
print(f"Score LOF moyen : {lof_scores.mean():.2f}")
\end{lstlisting}
\end{codebox}

\section{Conclusion}

LOF est l'outil de choix lorsque les anomalies sont contextuelles (ex: une température de 20°C est normale en été, mais anormale en hiver). Sa capacité à s'adapter à la densité locale le rend supérieur aux approches globales dans des datasets hétérogènes, au prix d'un coût de calcul plus élevé.

\end{document}
